\section*{Chapter \ref{ch:iv:behavioural-and-fit} --- Behavioural and Fit}

\begin{solution}{iq:iv:behavioural-and-fit:1}
A scoring answer is specific and personal: ``I like making decisions with incomplete information and finding out quickly whether I
was right. In my final year I ran a small prediction-market book in a university competition; I learned more from sizing positions
against people who knew more than me than from any course, and I want a job where that feedback is the job.'' It names what the
candidate likes in the work, shows it with an experience, and connects it to the role. It avoids money as the only reason, and
anything the interviewer could say about any job.
\iqlookfor{specificity, an example from the candidate's own life, and a link to the role.}
\end{solution}

\begin{solution}{iq:iv:behavioural-and-fit:2}
Tie the answer to the nature of the work: research questions measured over months with rigorous testing (against a market maker's
millisecond feedback and a bank's client-driven flow), models as the product, and, if true, a preference for depth over speed. Show
knowledge of what the firm type does day to day (One Quant Book~17, chapters 4 and~17), and avoid disparaging the alternatives.
\iqlookfor{an informed contrast between firm types that fits the candidate.}
\end{solution}

\begin{solution}{iq:iv:behavioural-and-fit:3}
In the \omterm{def:iv:behavioural-and-fit:star}{STAR structure}: the situation (two designs for a data pipeline), the task (you owned the ingestion part), the action (you
wrote down both options with their costs, ran a small benchmark or test that settled the factual part, and agreed on a decision
rule, deferring to the owner on what remained a judgement), the result (the chosen design, a number, and what you would do
differently). The interviewer scores that you argued with evidence, listened, and could disagree and commit.
\iqlookfor{a real disagreement resolved with evidence, told in the first person.}
\end{solution}

\begin{solution}{iq:iv:behavioural-and-fit:4}
The desk candidate names the loss with its size relative to limits, the view held and why, the moment the view was wrong, what was
done in the first hour (reduced risk, escalated, communicated), and what changed afterwards (a sizing rule, a stop, a check). The
graduate uses the nearest equivalent (a competition book, a failed project, a costly mistake in an internship) with the same
structure. Both separate decision from outcome: a good decision can lose, and saying so with the reasoning is more convincing than
presenting every loss as bad luck.
\iqlookfor{ownership, decision versus outcome, and a concrete change.}
\end{solution}

\begin{solution}{iq:iv:behavioural-and-fit:5}
$2.4/0.8 = 3$ standard deviations. If daily P\&L were normal, a loss that large would come about once in 741 days, about every three
years; with fat tails (a Student-$t$ with three degrees of freedom scaled to the same standard deviation) about once in 144 days,
every seven months. Tell the head of desk the number, the positions that produced it, whether it was within the book's stated risk
(a three-sigma day is rare but expected), whether anything in the book or the market has changed, and what you have already done:
the facts first, the explanation second, no defensiveness.
\iqlookfor{putting the loss in the book's own risk units, and a factual, prompt escalation.}
\end{solution}

\begin{solution}{iq:iv:behavioural-and-fit:6}
Choose a real bug: what it was, how it was found (ideally by you, or by monitoring you had set up), the impact with a number, the
immediate fix, and the systemic fix (a test, a check, an alert) so that the class of bug could not recur. Tell it without blaming
others and without minimising it; the interviewer scores the incident response and the learning, in the spirit of blameless
reviews (One Quant Book~15, chapter~28).
\iqlookfor{honest incident narration with a systemic fix.}
\end{solution}

\begin{solution}{iq:iv:behavioural-and-fit:7}
To a trader: what a new trader's first six months look like, and how risk is given to them. To a researcher: how ideas go from a
notebook to production, and how a research result is reviewed. To a head of desk: what distinguishes people who do well in their
first two years, and what the desk is trying to change this year. Avoid questions answered on the website, and questions about pay
before an offer.
\iqlookfor{questions that show thought about the role and gather real information.}
\end{solution}

\begin{solution}{iq:iv:behavioural-and-fit:8}
Do not trade the company's securities or anything related, do not tell anyone else, and do not recommend a trade to anyone: the
information may be inside information, and using it, disclosing it or recommending on it are prohibited (EU Market Abuse Regulation,
articles 8, 10 and~14, and their equivalents). Tell your compliance department at once, so that the name can go on the restricted
or watch list, and leave the conversation with your friend there; if you are in doubt whether it is inside information, treat it as
if it were. The interviewer is looking for the reflex, not for a legal analysis.
\iqlookfor{no trading, no disclosure, and prompt escalation to compliance.}
\end{solution}

\begin{solution}{iq:iv:behavioural-and-fit:9}
Verify the error first (rerun with the look-ahead removed and quantify the effect), then tell the committee's chair or your own
manager before the meeting, in writing, with the numbers, and keep trying to reach the colleague. Do not quietly fix and re-present
the colleague's work, and do not let the committee rely on a result you believe is wrong. The interviewer scores that the
candidate checks before accusing, escalates in time, and puts the decision's integrity above seniority.
\iqlookfor{verification, timely escalation and integrity over hierarchy.}
\end{solution}

\begin{solution}{iq:iv:behavioural-and-fit:10}
Report it at once to your head of desk and risk, and follow the firm's error procedure: the position is an error, not a trade, and its
profit belongs to the firm's error account, not to your P\&L. Reduce it as the procedure says. Then find how the size passed the
checks (a missing fat-finger limit) and propose the fix. Saying nothing because it made money is the answer that ends the interview.
\iqlookfor{immediate disclosure regardless of outcome, and a control fix.}
\end{solution}
